% Part: first-order-logic
% Chapter: axiomatic-deduction
% Section: deduction-theorem

% verification of properties of provability needed for maximally
% consistent sets in the completeness chapter.

\documentclass[../../../include/open-logic-section]{subfiles}

\begin{document}

\iftag{FOL}
      {\olfileid[tr]{fol}{axd}{ded}}
      {\olfileid[tr]{pl}{axd}{ded}}

\olsection{Tümdengelim Teoremi}

Gördüğümüz gibi aksiyomatik bir sistemde !!{derivation}s vermek zahmetlidir ve
!!{derivation}s bulmak zor olabilir. Uzun !!{formula}s listelerini gerçekten
yazmak yerine, birkaç basit sonuçtan yararlanarak bu tür !!{derivation}s
bulunduğunu savunmak genellikle daha kolaydır. Böyle üç sonuç gösterdik:
\olref[ptn]{prop:reflexivity}, $!A \in \Gamma$ olduğunu bildiğimizde her zaman
$\Gamma \Proves !A$ diyebileceğimizi söyler. \olref[ptn]{prop:monotonicity},
$\Gamma \Proves !A$ ise $\Gamma \cup \{!B\} \Proves !A$ da olduğunu söyler.
\olref[ptn]{prop:transitivity} ise $\Gamma \Proves !A$ ve $!A \Proves !B$
olduğunda $\Gamma \Proves !B$ sonucunu verir. İşte bir başka basit sonuç,
modus ponensin ``meta'' sürümü:

\begin{prop}
  \ollabel{prop:mp} $\Gamma \Proves !A$ ve $\Gamma \Proves !A \lif !B$ ise
  $\Gamma \Proves !B$.
\end{prop}

\begin{proof}
  $\{!A, !A \lif !B\} \Proves !B$ olduğunu şu türetim gösterir:
  \begin{derivation}
    1. & $!A$ & Vars.\\
    2. & $!A \lif !B$ & Vars.\\
    3. & $!B$ & 1, 2, MP
  \end{derivation}
  \olref[ptn]{prop:transitivity} ile $\Gamma \Proves !B$ elde edilir.
\end{proof}

Bu bağlamda kullanacağımız en önemli sonuç tümdengelim teoremidir:

\begin{thm}[Tümdengelim Teoremi]
\ollabel{thm:deduction-thm} $\Gamma \cup \{!A\} \Proves !B$ ancak ve ancak
$\Gamma \Proves !A \lif !B$ ise geçerlidir.
\end{thm}

\begin{proof}
``Eğer'' yönü doğrudandır. $\Gamma \Proves !A \lif !B$ ise
\olref[ptn]{prop:monotonicity} uyarınca $\Gamma \cup \{!A\} \Proves
!A \lif !B$. Ayrıca \olref[ptn]{prop:reflexivity} ile
$\Gamma \cup \{!A\} \Proves !A$. Dolayısıyla \olref{prop:mp} uyarınca
$\Gamma \cup \{!A\} \Proves !B$.

``Ancak'' yönünde, $\Gamma \cup \{!A\}$'dan $!B$ için verilen
!!{derivation} ele alınır ve uzunluğuna göre tümevarım yapılır.

Tümevarımın başlangıcında iddiayı uzunluğu~$1$ olan her !!{derivation} için
ispatlarız. $\Gamma \cup \{!A\}$'dan $!B$ için uzunluğu~$1$ olan
!!^a{derivation}, tek başına $!B$'den oluşur; doğruysa $!B$ ya
$\Gamma \cup \{!A\}$ içindedir ya da bir aksiyomdur. $!B \in \Gamma$ ise
veya $!B$ bir aksiyomsa $\Gamma \Proves !B$. Ayrıca
\olref[prp]{ax:lif1} uyarınca $\Gamma \Proves !B \lif (!A \lif !B)$;
\olref{prop:mp} de $\Gamma \Proves !A \lif !B$ verir. $!B \in \{!A\}$
ise $!B$ ile $!A$ aynı formül olduğundan $\Gamma \Proves !A \lif !B$;
son !!{formula}~$!A \lif !B$,
$!A \lif !A$ ile aynıdır; bunu \olref[pro]{ex:identity} içinde
türettik (\emph{!!{derive}}).

Tümevarım adımında, $\Gamma \cup \{!A\}$'dan $!B$ için bir
!!a{derivation} ele alalım ve bunun modus ponens ile gerekçelendirilen $!B$
adımıyla bittiğini varsayalım. (Son adım modus ponensle gerekçelendirilmemişse,
$!B$ bir aksiyom, $\Gamma$'nın elemanı ya da $!A$ olduğunda başlangıç
adımındaki akıl yürütme geçerlidir; FOL'daki niceleyici kuralı durumu sonraki
bölümde ayrıca ele alınır.) O zaman !!{derivation}
içindeki önceki adımlar, bir
!!{formula}~$!C$ için $!C \lif !B$ ve $!C$'dir; yani
$\Gamma \cup \{!A\} \Proves !C \lif !B$ ve
$\Gamma \cup \{!A\} \Proves !C$. İlgili !!{derivation}s daha kısa
olduğundan tümevarım varsayımı uygulanır. Böylece
\begin{align*}
  & \Gamma \Proves !A \lif (!C \lif !B); \\
  & \Gamma \Proves !A \lif !C
\end{align*}
elde ederiz. Ayrıca \olref[prp]{ax:lif2} ile
\[
\Gamma \Proves (!A \lif (!C \lif !B)) \lif
((!A\lif !C) \lif (!A \lif !B)).
\]
\olref{prop:mp}'nin iki uygulaması, istendiği gibi
$\Gamma \Proves !A \lif !B$ verir.
\end{proof}

\olref[prp]{ax:lif1} ile \olref[prp]{ax:lif2}'nin, Tümdengelim Teoremi tam
olarak geçerli olsun diye seçildiğine dikkat edin.

Aşağıda !!{derivability} hakkında bazı yararlı olgular verilmiştir; bunları
alıştırma olarak bırakıyoruz.

\begin{prop}
  \ollabel{prop:derivfacts}
\begin{enumerate}
\item $\Proves (!A \lif !B) \lif ((!B \lif !C)
  \lif (!A \lif !C))$; \ollabel{derivfacts:a}
\item $\Gamma \cup \{ \lnot !A\} \Proves \lnot !B$ ise
  $\Gamma \cup \{ !B\} \Proves !A$ (Karşıt konum); \ollabel{derivfacts:b}
\item $\{ !A, \lnot!A\} \Proves !B$ (Ex Falso Quodlibet, Patlama);
  \ollabel{derivfacts:c}
\item $\{ \lnot\lnot!A\} \Proves !A$ (Çifte Değillemenin Giderilmesi);
  \ollabel{derivfacts:d}
\item $\Gamma \Proves \lnot\lnot!A$ ise $\Gamma \Proves !A$;
  \ollabel{derivfacts:e}
\end{enumerate}
\end{prop}

\tagprob{FOL}
\begin{prob}
  \olref[fol][axd][ded]{prop:derivfacts} önermesini ispatlayın.
\end{prob}
\tagendprob

\tagprob{notFOL}
\begin{prob}
  \olref[pl][axd][ded]{prop:derivfacts} önermesini ispatlayın.
\end{prob}
\tagendprob

\end{document}
